% !TeX root = ../教材题库_全部条目_转写工作稿.tex
% Source: XJ-B1-C05-TH-P165-Y1147
% 第5章 / 5.2 任意角的三角函数 / 教材 p.165、166 / PDF 第172、173页
% OCR未校对：公式与文字均可能有误；未完成数学TeX转写；禁止直接发给学生
% 原题及相邻原解截图（核对用）：../题目原图/XJ-B1-C05-TH-P165-Y1147-01.webp, ../题目原图/XJ-B1-C05-TH-P165-Y1147-02.webp
% 本文件只在转写完成后作为题干使用；OCR 条目尚不是可靠题干。
\begin{bankitem}{p.165 思考（含紧邻说明）}{第5章 / 5.2 任意角的三角函数 / 教材 p.165、166 / PDF 第172、173页}{XJ-B1-C05-TH-P165-Y1147}{OCR未校对：公式与文字均可能有误；未完成数学TeX转写；禁止直接发给学生}
\begin{ocrdraft}
[教材p.165 / PDF第172页]
\par
？思考我们不难发现，任意角的正弦值、余弦值与正切值，既有可能是正
\par
数，也有可能是负数，还可能为0．它们的符号与什么有关？
\par
设P(x，3，）是角a的终边上不同于原点的任意一点，记厂一十重，则sina=
\par
由7．>0可知，sin的正负与角的终边上点的纵坐标的符号相同．
\par
于是，对于当且仅当角仅的终边在第一、二象限或3，轴正半轴上时，
\par
sina>0；当且仅当角的终边在第三、四象限或轴负半轴上时，sina<0.
\par
用类似方法可以得到：
\par
对于cosa=
\par
，当且仅当角的终边在第一、四象限或工轴正半轴上时，>
\par
0；当且仅当角a的终边在第二
\par
三象限或轴负半轴上时，<0．
\par
[教材p.166 / PDF第173页]
\par
对于=
\par
，当且仅当角a的终边在第一
\par
。象限时，tana>0；当且仅当角
\par
a的终边在第二、四象限时，tana<0.
\par
以上结论可用图5，2一5直观表示．
\par
SIn仅
\par
COS
\par
tan“
\par
0
\par
0
\par
0
\par
（2）
\par
（3）
\par
图5．2一5
\end{ocrdraft}
\end{bankitem}
